\section{A split \(\mathbb Q(\sqrt{-7})\) sixfold source from \(G_{24}\)}
\label{sec:g24-new}

Let \(K=\mathbb Q(\tau)\), where
\(\tau=-(1+\sqrt{-7})/2\), so \(\tau^2+\tau+2=0\) and
\(\mathcal O_K=\mathbb Z[\tau]\).  We construct a dominant family of split
Weil sixfolds, supply its determinant classes, and give an explicit effective
cycle bound for every fixed split integral datum over this field.

\paragraph{Provenance.}
The initial three reflection matrices are the specialization \(p=2\) of
Deraux's published matrices \cite[equation (13), author manuscript p.~24;
\S3, pp.~4--6]{G24:Deraux}.  The six roots and nine boundary intervals below
were found in this run by an exact search, beginning with those three matrices.
They are not copied from Severn or rStation.  The unrefereed rStation papers
58 and 75 suggested looking at \(G_{24}\), but no statement or numerical
witness from them is a premise.  The fixed certificate
\texttt{g24/g24\_reconstruct.py} reconstructs every matrix from the displayed
roots.  No numerical approximation, claimed period-rank routine, or dimension
inequality is used to prove dominance.

\subsection*{D-G24-01. The finite group, integral action, and branch data}

Put
\[
 H=\begin{pmatrix}
 2&-\tau&\bar\tau\\
 -\bar\tau&2&-\tau\\
 \tau&-\bar\tau&2
 \end{pmatrix},\qquad
 \langle x,y\rangle=\bar x^{\,t}Hy.
\]
Its leading principal minors are \(2,2,1\), so it is positive definite.
For a nonzero root \(v\), define
\[
 s_v=1-\frac{2v\bar v^{\,t}H}{\langle v,v\rangle}.
\]
The coordinate roots give Deraux's three involutions
\[
 R_1=\begin{pmatrix}-1&\tau&-\bar\tau\\0&1&0\\0&0&1\end{pmatrix},\quad
 R_2=\begin{pmatrix}1&0&0\\\bar\tau&-1&\tau\\0&0&1\end{pmatrix},\quad
 R_3=\begin{pmatrix}1&0&0\\0&1&0\\-\tau&\bar\tau&-1\end{pmatrix}.
\]
Their group is the rank-three reflection group \(G_{24}\), of order
\(336\).  In addition to the published identification, a direct enumeration
over \(K\) gives exactly \(336\) distinct matrices, all preserving \(H\).
Their orbit on the first coordinate vector has \(42\) vectors of norm two,
or \(21\) root lines.  These are explicit finite equality checks, not a
recognition based only on the order.

Take the following six norm-two roots:
\[
\begin{array}{c|rrr}
 i&\multicolumn{3}{c}{v_i^t}\\\hline
 1&0&-1&1+\tau\\
 2&-1&0&\tau\\
 3&-2&\tau&1+\tau\\
 4&-1-\tau&-2&\tau\\
 5&-1&0&0\\
 6&-1&-1&\tau
\end{array}
\]
Set \(s_i=s_{v_i}\) and
\[
                (t_1,\ldots,t_{12})=(s_1,s_1,\ldots,s_6,s_6).
\]
The exact enumeration verifies that the \(s_i\) generate the same \(336\)
matrices as \(R_1,R_2,R_3\).  Their root span has rank three and their
nonorthogonality graph is connected.  The orthogonal-summand argument used
in D-G34-01 therefore proves absolute irreducibility over \(\mathbb C\).
All matrices lie in \(K\), and
\[
                  \operatorname{tr}(R_1R_2R_3)=1+\tau.
\]
Thus the character field is exactly \(K\).  The displayed realization has
Schur index one, so the rational simple block is \(M_3(K)\).

The doubled tuple has product one and generates the group.  Riemann
existence gives a connected \(G_{24}\)-cover \(C\to\mathbb P^1\), with
twelve branch points of inertia order two.  Its Hurwitz component is a
finite cover of \(M_{0,12}\).  Riemann--Hurwitz gives
\[
 2g(C)-2=336\bigl(-2+12(1-1/2)\bigr)=4\cdot336,
 \qquad g(C)=673.
\]
For the natural character \(\sigma\) and its conjugate, each local
eigenvalue list is \((1,1,-1)\).  Chevalley--Weil
\cite[Theorem 2.10]{G34:FGP} gives holomorphic multiplicity
\[
                         -3+12/2=3
\]
at each embedding.  A rational primitive idempotent of \(M_3(K)\) therefore
cuts out an abelian sixfold of signature \((3,3)\).

We need its integral structure, not just its rational isogeny class.  Let
\(G=336\), let \(\chi\) be the natural character, and put
\[
 \varepsilon=\frac3G\sum_{a\in G_{24}}
       \operatorname{Tr}_{K/\mathbb Q}\bigl(\chi(a^{-1})\bigr)a,
 \qquad e=\varepsilon\frac{1-s_1}{2}.
\]
The character traces are integers.  Thus \(2Ge\in\mathbb Z[G_{24}]\).
The exact coefficient calculation also verifies that \(Ge\) itself has
denominator two; it would be incorrect to reuse the even-dimensional
\(G_{34}\) projector denominator.  The element \(e\) is a primitive
Rosati-symmetric rational idempotent.  Let \(A\subset J(C)\) be its image,
with the restriction \(\theta_A\) of the principal polarization.

There is an integral \(\mathcal O_K\)-action on this particular subvariety.
Let \(g_0=R_1R_2R_3\).  Direct multiplication gives
\[
 g_0=\begin{pmatrix}
 1+\tau&1&1\\1&1&-\tau\\-\tau&-1-\tau&-1
 \end{pmatrix},\qquad
 \langle v_1,g_0v_1\rangle=\tau,
 \qquad (1-s_1)g_0v_1=\tau v_1.
\]
The integral endomorphism \((1-s_1)g_0\) of \(J(C)\) preserves \(A\).
Indeed \(g_0\) preserves the central \(\varepsilon\)-block and \(1-s_1\)
maps that block to \(A\).  On the natural root line its restriction is
\(\tau\), and on the conjugate line it is \(\bar\tau\).  Its restriction
to \(A\) thus supplies \(\tau\in\operatorname{End}(A)\), proving
\(\mathcal O_K\)-stability of \(H_1(A,\mathbb Z)\) without saturation.

Write \(\Lambda=H_1(J(C),\mathbb Z)\) and
\(\Lambda_A=\Lambda\cap e\Lambda_{\mathbb Q}\).  The alternating
dual is \(\Lambda_A^\#=e\Lambda\): the primitive inclusion of
\(\Lambda_A\) in the unimodular lattice \(\Lambda\) lets every integral
functional extend, and self-adjointness of \(e\) identifies its representing
vector with its projection.  Hence
\[
 2G\Lambda_A^\#\subseteq\Lambda_A.
\]
Every polarization elementary divisor divides \(2G=672\).  Therefore
\begin{equation}\label{eq:g24-source-delta}
 \Delta_A\mid(2G)^6,
 \qquad \int_A\theta_A^6=6!\Delta_A\le6!(2G)^6.
\end{equation}

\subsection*{D-G24-02. Nine exact boundary logarithms}

Use left cocycles \(u(gh)=u(g)+g u(h)\), and write
\[
 u_i=\lambda_i v_{\lceil i/2\rceil},\qquad Q_i=t_1\cdots t_{i-1}.
\]
The relation is \(\sum_i Q_i u_i=0\); the coboundaries are
\(u_i=(1-t_i)a\).  The coefficient matrices have ranks three and three,
their product is zero, and the quotient dimension is \(12-3-3=6\).

Let \(b_1,b_2,b_3\) be the equal-pair cocycles supported respectively at
\((1,2),(3,4),(5,6)\), with coefficient one at each member of the pair.
The remaining representatives are supported at \((1,3,5,7,9,11)\), with
rows
\[
\begin{aligned}
 b_4:\;&(-1-\tau/2,\ 1,\ -1-\tau/2,\ 1,\ 0,\ 0),\\
 b_5:\;&(1-\tau/2,\ 1+\tau,\ -1-\tau/2,\ 0,\ 1,\ 0),\\
 b_6:\;&(-1-\tau,\ 1,\ -1,\ 0,\ 0,\ 1).
\end{aligned}
\]
All six satisfy the relation.  Adjoining the three coboundary columns gives
rank nine, so they are a quotient basis.  Set
\(I=\langle b_1,b_2,b_3\rangle_K\).

As an auxiliary algebraic certificate, set
\[
 S(u,v)=\sum_i\left\langle Q_i^{-1}\sum_{j<i}Q_j u_j+u_i/2,v_i\right\rangle.
\]
The reconstruction checks \(S^*=-S\), \(\operatorname{rank}S=6\),
and a zero upper-left \(3\)-by-\(3\) block.  Its upper-right block is
\[
 B=\begin{pmatrix}
 -1-\tau/2&1-\tau/2&-1-\tau\\
 1&1+\tau&1\\
 -1-\tau/2&-1-\tau/2&-1
 \end{pmatrix}.
\]
The proof uses this form only for its explicitly verified matrix properties;
no unproved conversion between cup-product conventions is required.

Take the nine intervals
\[
 [1,2],\ [3,4],\ [5,6],\ [7,8],\ [9,10],\ [11,12],
 \ [1,4],\ [5,8],\ [1,8].
\]
They are pairwise nested or disjoint and define distinct stable partitions.
Thus they are a maximal compatible collection of nine boundary circles on
the twelve-punctured sphere.  Their products \(c=t_a\cdots t_b\) are all
the identity.  In particular every boundary twist already returns the
monodromy tuple, without a nontrivial return power.

The twist formula is
\[
 u(c)=\sum_{i=a}^b(t_a\cdots t_{i-1})u_i,\qquad
 u_i'=u_i+(1-t_i)u(c)\quad(a\le i\le b),
\]
with unchanged outside coordinates.  It follows by applying the cocycle
identity to conjugation by the interval boundary word.  Reducing these
transformed representatives modulo coboundaries gives
\[
 T_j=1+N_j,\qquad N_j=\begin{pmatrix}0&X_j\\0&0\end{pmatrix},
 \qquad N_iN_j=0,\qquad T_j^*ST_j=S.
\]
The ranks of the \(N_j\), in the displayed interval order, are
\[
                         1,1,1,1,1,1,2,2,2.
\]
All nine \(3\)-by-\(3\) blocks and the full matrix \(S\) are printed
in the associated exact matrix appendix.

For each \(j\), let \(Q_j'\) be minus the lower-right block of \(SN_j\).
It is Hermitian positive semidefinite.  Here is the complete principal-minor
certificate, with columns ordered as
\(1,2,3,12,13,23,123\):
\[
\begin{array}{c|rrrrrrr}
 j&1&2&3&12&13&23&123\\\hline
1&2&4&4&0&0&0&0\\
2&2&4&2&0&0&0&0\\
3&2&2&2&0&0&0&0\\
4&2&0&0&0&0&0&0\\
5&0&2&0&0&0&0&0\\
6&0&0&2&0&0&0&0\\
7&3/2&5/2&2&2&1&1&0\\
8&3/2&2&2&2&2&0&0\\
9&0&2&2&0&0&2&0
\end{array}
\]
The leading principal minors of \(\sum_jQ_j'\) are
\[
                              11,\quad143,\quad732.
\]
Thus every positive weighted sum of the \(Q_j'\) is positive definite.
Since the lower-left block of \(S\) is invertible, every positive weighted
sum of the \(N_j\) has rank three and has image and kernel both \(I\).

Flatten the nine \(X_j\) row by row.  The resulting \(9\)-by-\(9\)
matrix over \(K\) has determinant
\begin{equation}\label{eq:g24-period-minor}
             -308-126\tau\ne0,
 \qquad N_{K/\mathbb Q}(-308-126\tau)=87808=2^8\,7^3.
\end{equation}
This proves complex linear independence under either embedding of \(K\),
not merely rational independence of a doubled list of coefficients.

\subsection*{D-G24-03. The analytic rank bridge and splitness}

We give the analytic step explicitly.  The nine circles determine a
zero-dimensional boundary stratum of \(\overline M_{0,12}\), with nine
independent plumbing coordinates \(z_1,\ldots,z_9\).  After a finite cover
of the Hurwitz space marking the finite monodromy, and a coordinatewise
Kummer base change, the associated polarized weight-one variation has
unipotent monodromy logarithms
\[
 M_j=\epsilon c_jN_j,\qquad c_j\in\mathbb Z_{>0},
 \qquad\epsilon\in\{1,-1\}\text{ independent of }j.
\]
The common sign reflects the global choice of pullback versus transport
and orientation; it cannot be chosen independently at different nodes.
The geometric parabolic action is the cocycle action by
\cite[Lemma 2.2 and Theorem 2.5]{G34:DWmon}.  Finite residual central
ambiguities are removed by the finite base change.  Since \(N_j^2=0\),
the logarithm is exactly the indicated scalar multiple of \(N_j\).

Put \(z_j=e^{2\pi i w_j}\).  The nilpotent orbit theorem
\cite[Theorem 4.12, p.~233]{G34:Schmid} gives
\[
 F(w)=\exp\!\left(\sum_jw_jM_j\right)\Psi(z),
 \qquad\Psi\text{ holomorphic at }0,
 \qquad F_0=\Psi(0).
\]
Along the diagonal, \(M=\sum_jM_j\) has rank three over \(K\), and
\(M^2=0\), \(\operatorname{im}M=\ker M=I\).  On the underlying
twelve-dimensional rational space the monodromy weight filtration is
\[
                    W_0=W_1=I_{\mathbb Q},\qquad W_2=V_{\mathbb Q}.
\]
Schmid's limiting mixed Hodge structure theorem
\cite[Theorem 6.16, p.~255]{G34:Schmid} applies.  The original weight-one
filtration has \(F^0=V_{\mathbb C}\), \(F^2=0\); these conditions remain
true at \(F_0\).  On \(\operatorname{Gr}_0^W\), both Hodge indices are
therefore nonnegative and sum to zero.  It has type \((0,0)\), so
\[
                            F_0^1\cap I_{\mathbb C}=0.
\]
In either embedding of \(K\), \(F_0^1\) is a three-dimensional complement
to \(I\) in a six-dimensional space.

The tangent space of \(\operatorname{Gr}(3,6)\) there is
\(\operatorname{Hom}(F_0^1,I)\).  After projection along the complement,
the tangent vectors \(M_jF_0\) are represented by \(\epsilon c_jX_j\),
up to the same invertible basis changes.  Equation
\eqref{eq:g24-period-minor} makes all nine independent.

Translate a period derivative back by
\(\exp(-\sum_jw_jM_j)\).  In any local Grassmannian chart it is
\[
              M_j\Psi(z)+2\pi i z_j\frac{\partial\Psi}{\partial z_j}(z).
\]
As every \(z_j\) tends to zero through nonzero values, its tangent value
tends to \(M_jF_0\).  The determinant of these nine translated tangent
vectors thus has a nonzero limit.  It is nonzero at sufficiently nearby
smooth points.  The actual period map has rank nine there; the unitary
period domain of signature \((3,3)\) has dimension nine.  The period
image consequently contains a nonempty analytic open subset.

Finally \(M\), as a geometric logarithm, is skew-adjoint for the genuine
alternating polarization.  Thus
\[
 E(Mu,Mv)=-E(u,M^2v)=0.
\]
Its image \(I\) is a three-dimensional isotropic \(K\)-subspace.  Hence
the Hermitian form is split.  This conclusion uses the actual polarized
variation; it does not identify the auxiliary matrix \(S\) with the
geometric cup product by fiat.

\subsection*{D-G24-04. Algebraic determinant classes and transfer}

The supplier and transfer proofs D-G34-04--06 apply to this cover.  We
specify all their hypotheses and the dimensional changes to avoid hiding
an additional theorem premise.

For a subgroup \(H'\le G_{24}\) and a linear character \(\chi\) of order
\(m\ge3\), the intermediate map
\(C/\ker\chi\to C/H'\) is cyclic.  Every nonzero inertia image is the
unique element of order two.  If \(m\) is odd the map is unramified; if
\(m\) is even, all nonzero branch exponents equal \(m/2\).  The abelianized
surface-group relation makes their number even, so they pair with their
negatives.  For base genus \(q\) and \(s\) nonzero branch points, every
primitive eigenspace has even rank \(h=2q-2+s\) and signature
\((h/2,h/2)\).  Schoen's cyclic-product theorem
\cite[Theorem 2.0 and Corollary 3.1]{G34:Schoen}, with the faithful-character
localization proof of D-G34-04 replacing the defective prime-quotient step,
supplies rational algebraic determinant classes.  That repair is independent
of \(G_{34}\), the field of the selected target character, and the target
dimension.  It uses only paired nonzero branch exponents.  Its original
center and exceptional-divisor vanishings both apply here.

The rational primitive cyclotomic determinant space has rank one over its
cyclotomic field.  A nonzero rational element has a nonzero component at
every primitive embedding, so the suppliers cover all required characters.
Let \(U=\operatorname{Cent}_{\operatorname{Sp}(H^1(C))}(G_{24})\), including
its disconnected orthogonal factors.  Brauer induction gives an integral
relation for \(\sigma\), hence for its determinant character in
\[
                       X^*(U_{\mathbb C})=\mathbb Z^a\oplus(\mathbb Z/2)^b.
\]
The orders-one-and-two terms vanish: free coordinates cancel under complex
conjugation, and quaternionic restriction multiplicities are even by the
nondegenerate alternating form on a real-character eigenspace.  Zero
cohomology terms also vanish.  The remaining terms are exactly the supplied
cyclic determinants.  Replacing a character by its inverse handles negative
coefficients in both the free and parity coordinates.

Choose one primitive determinant line for each term and its complementary
exterior line.  The integral character relation makes their product with
the desired line \(\det M_\sigma\subset H^6(A,\mathbb C)\) \(U\)-invariant.
All factors are algebraic isogeny factors of \(J(C)\) via quotient pullbacks
and rational group-algebra projectors.  A \(U\)-equivariant splitting lifts
the tensor to a power of \(J(C)\).  The Lefschetz group of \(J(C)\) is
contained in \(U\), so \cite[Theorem 3.2]{G34:Milne} makes this tensor a
divisor-generated algebraic class.  Contracting the complementary lines
against the rational suppliers and pushing forward isolates the desired
primitive determinant line with a nonzero coefficient.  Repeat for the
conjugate line and descend from complex to rational coefficients.  The
whole rational Weil plane is algebraic.

For the quantitative argument use the middle-degree version on the curve
powers: a primitive determinant in \(H^1(C_i)^{\otimes h_i}\) pairs with
the inverse-character determinant in the same degree \(h_i\).  If
\(\ell=\sum_i|n_i|h_i\), the invariant contracting tensor on the product
of these curve powers with \(A\) has degree \(\ell+6\).  Thus its divisor
degree is \(\ell/2+3\).  This is the sixfold replacement of
\(\ell/2+4\) in the eightfold construction, and involves no inverse
Lefschetz operation of uncontrolled degree.

\subsection*{D-G24-05. An explicit detector bound}

Define the following positive integers:
\[
\begin{gathered}
 G=336,\qquad g=673,\qquad M=2G=672,\\
 R=(2G+1)G!G^G,\qquad L=2gR,\qquad
 S_0=G^{2g}(g!)^2,\qquad S_B=(L/2)!S_0^R,\\
 k=L/2+3,\qquad
 B_{0,7}=(40G)^k\binom{k}{3}S_B\,6!M^6.
\end{gathered}
\]
For a fixed split \(\mathcal O_K\)-stable Weil datum with polarization
elementary-divisor product \(\Delta\), put
\begin{equation}\label{eq:g24-detector-bound}
 N_7=4M^6\Delta,\qquad
 \boxed{D_{6,7}(\Delta)=N_7^6B_{0,7}.}
\end{equation}
Every fibre of its connected fixed-data family has an effective integral
three-cycle with nonzero Weil projection and degree at most
\(D_{6,7}(\Delta)\).

\begin{proof}
We adapt the fully checked arithmetic and intersection arguments SQ3--SQ7
of \cite{G34:Hecke}; their original \(G_{34}\) geometric premises have
already been replaced by D-G24-01--04.  The integral relation has coefficient
mass at most \(R\): the character lattice has at most \(G\) rows, all
entries have absolute value at most \(G\), and appending the columns
\(2e_j\) treats the orthogonal parities.  The integer-lattice determinant
argument of SQ3 gives \((2G+1)G!G^G\), independent of the target dimension.
The cyclic primitive projectors are cleared by \(G\), since their
coefficients are Ramanujan sums divided by subgroup orders.  Supplier ranks
are at most \(g\), and the integral orbit-sum cycle mass is bounded by
\(S_0\) as in SQ4.  The product supplier on the curve product has degree
mass at most \(S_B\), with \(\ell\le gR\le L\), as in SQ5.

The full-centralizer invariant algebra is generated by degree-two tensors
coming from \(\mathbb Q[G_{24}]\), with disconnected orthogonal factors
included.  The effective divisor spanning and moving argument of SQ6
therefore applies.  Its globally generated divisors are bounded in the nef
order by \(40G(H_B+\theta_A)\).  The bound \(40G\) depends on the finite
cover degree and on the curve-to-Jacobian pullbacks, not on the representation
dimension.  Intersect the signed product supplier times \(A\) with a
divisor monomial of degree \(c=\ell/2+3\), then push the resulting
three-cycles to \(A\).  Their total degree mass is at most
\[
\begin{aligned}
 &(40G)^c\int_{|Y_B|\times A}(H_B+\theta_A)^c\theta_A^3\\
 &\qquad=(40G)^c\binom c3
       \operatorname{mass}_{H_B}(Y_B)\int_A\theta_A^6
 \le B_{0,7}.
\end{aligned}
\]
The invariant contraction is in the complex span of these divisor monomials
and has nonzero Weil projection.  Therefore at least one of the finite
collection of effective integral output cycles has nonzero Weil projection.
In fact their projected classes span both Weil eigenlines.  No bound on
the scalar coefficients of a complex divisor-polynomial presentation is
needed for this extraction.

It remains to compare integral split lattices.  For the alternating form
\(E\), define
\[
                  h(x,y)=E(x,\tau y)-\bar\tau E(x,y).
\]
This is an integral \(\mathcal O_K\)-valued Hermitian form, and the
coefficient of \(\tau\) recovers \(E\).  In an \(\mathcal O_K\)-basis,
\(\det E=(\det h)^2\); in signature \((3,3)\), \(\det h=-\Delta\).
The ring \(\mathcal O_K\) is a PID.  For completeness, it is norm
Euclidean: choose the \(\tau\)-coefficient of an approximation so that the
imaginary error has absolute value at most \(\sqrt7/4\), then choose the
integer coefficient so that the real error has absolute value at most
\(1/2\).  The error norm is at most \(7/16+1/4=11/16<1\).

A primitive maximal isotropic submodule therefore extends to an integral
basis with matrix
\[
                 \begin{pmatrix}0&B\\B^*&D\end{pmatrix},
 \qquad b=\det B,\qquad N_{K/\mathbb Q}(b)=\Delta.
\]
Writing \(X=(B^*)^{-1}D/2\), the substitutions \(w'=w-vX\) and
\(u=w'B^{-1}\) make a standard hyperbolic basis.  The forward and inverse
coordinate matrices have denominators dividing \(2\Delta\): the only
denominators are \(2\bar b,b,2b\bar b\).  Applying this to source and
target gives a comparison denominator dividing \(4\Delta_A\Delta\),
which is cleared by \(N_7=4M^6\Delta\) in view of
\eqref{eq:g24-source-delta}.  Multiplying the rational isometry by \(N_7\)
gives an isogeny with
\[
                  \phi^*\theta_{\mathrm{target}}=N_7^2\theta_A.
\]
Its pushforward multiplies the degree of a three-cycle by \(N_7^6\),
and preserves nonzero Weil projection.  D-G24-03 gives an analytic open
period image, so \eqref{eq:g24-detector-bound} holds on an analytic open
subset of the fixed target family.  The project's accepted bounded-cycle
closedness and transport theorem extends that same degree bound to every
fibre of its connected component.
\end{proof}

For a fixed nonmaximal order \(\mathcal O_f=\mathbb Z+f\mathcal O_K\),
put \(\Lambda'=f\mathcal O_K\Lambda\).  Then
\[
 f\Lambda\subset\Lambda'\subset\Lambda,\qquad
 I=[\Lambda:\Lambda']\le f^{12},\qquad \Delta'=I\Delta.
\]
The identity isogeny from the torus with lattice \(\Lambda'\) to the
original torus has pullback polarization of volume \(\Delta'\).  Pushing
an effective three-cycle forward preserves its degree for this pair of
polarizations.  Thus
\[
                     \boxed{\delta_W\le D_{6,7}(f^{12}\Delta).}
\]
This is a fixed-order assertion, with the conductor dependence displayed.
Applying the project's accepted Weil degree-to-division theorem gives the
corresponding explicit bounded-division estimate for every fixed rational
Weil class.  In particular the principal balanced \(E^6\) datum over
\(K=\mathbb Q(\sqrt{-7})\) is covered.  The argument establishes the
split sixfold component; it does not assert an all-determinant sixfold
theorem over this field.
